\capitulo[Nome da Autora]{Título do segundo capítulo}

\begin{resumo}
Resumo do segundo capítulo.
\end{resumo}
\palavraschave{termo; termo; termo.}

\section{Seção}
Cada capítulo abre em página ímpar. A hemogasometria\index{Hemogasometria} também aparece aqui, e o índice remissivo reúne as
páginas automaticamente. Em cães\index{Cães} e gatos\index{Gatos}, os valores de referência diferem.

\section{Considerações finais}
Síntese e aplicações práticas do conteúdo.

\begin{referencias}
\item SOBRENOME, Nome. \textbf{Título}. Cidade: Editora, 2022.
\end{referencias}
